\section{Implementation Skills}

\begin{purposebox}
\textbf{Purpose:} Demonstrate readiness to implement the tracking control NN in MATLAB, Python, and C++/CUDA. Not about mastery of each language, but about knowing the right tools and libraries for control system simulation.

\textbf{When complete:} You can set up a simulation environment in each language and implement basic control system operations.
\end{purposebox}

%──────────────────────────────────────────────────
\subsection{MATLAB Implementation}
%──────────────────────────────────────────────────

\subsubsection*{Concept Check}

\question{Core Functions}
For each MATLAB function, describe what it does and when you would use it in this project:

\medskip
\begin{tabularx}{\textwidth}{l X X}
\toprule
\textbf{Function} & \textbf{Purpose} & \textbf{When to Use} \\
\midrule
\texttt{ss()}   & & \\[0.8em]
\texttt{tf()}   & & \\[0.8em]
\texttt{step()} & & \\[0.8em]
\texttt{bode()} & & \\
\bottomrule
\end{tabularx}

\begin{answerbox}\end{answerbox}

\question{Simulink}
What is Simulink and how does it relate to the functional block diagram from Step~4? Would you implement the block diagram directly in Simulink or in MATLAB script? What are the trade-offs?

\begin{answerbox}\end{answerbox}

\question{MATLAB Simulation Skeleton}
Write pseudocode (or actual MATLAB code) for the basic simulation loop of the tracking control NN:
\begin{itemize}[nosep]
  \item Initialize state $x(0)$, reference trajectory $r(k)$, parameters ($\tau, \beta, f, g$)
  \item Loop over $k = 0$ to $N$: compute $u(k)$, update $x(k+1)$, store results
  \item Plot $x(k)$ vs.\ $r(k)$
\end{itemize}

\begin{codeanswer}\end{codeanswer}

%──────────────────────────────────────────────────
\subsection{Python Implementation}
%──────────────────────────────────────────────────

\subsubsection*{Concept Check}

\question{Library Mapping}
Map MATLAB functionality to Python libraries:

\medskip
\begin{tabularx}{\textwidth}{l l X}
\toprule
\textbf{MATLAB} & \textbf{Python Equivalent} & \textbf{Library} \\
\midrule
Matrix operations       & & \\[0.6em]
Control System Toolbox  & & \\[0.6em]
ODE / system simulation & & \\[0.6em]
Plotting                & & \\
\bottomrule
\end{tabularx}

\begin{answerbox}\end{answerbox}

\question{NumPy for the System Model}
Write Python code (using NumPy) to compute one time step of the system model:
\[
  x(k+1) = x(k) + \tau\bigl(f(x(k)) + g(x(k))\,u(k)\bigr)
\]
Assume \texttt{x} is an $n$-dimensional array, \texttt{u} is an $m$-dimensional array, and \texttt{f()}, \texttt{g()} are functions.

\begin{codeanswer}\end{codeanswer}

\question{\texttt{python-control} Library}
What does the \texttt{python-control} library provide? How would you use it for this project? Give an example function call.

\begin{answerbox}\end{answerbox}

%──────────────────────────────────────────────────
\subsection{C++ and CUDA}
%──────────────────────────────────────────────────

\subsubsection*{Concept Check}

\question{Why C++?}
Why would you implement the tracking control NN in C++ when MATLAB and Python are easier? What practical applications require C++ performance?

\begin{answerbox}\end{answerbox}

\question{Eigen Library}
The Eigen library provides matrix operations in C++. Write pseudocode showing how you would represent the state vector $x(k)$ and compute one system update step using Eigen-style operations.

\begin{codeanswer}\end{codeanswer}

\question{CUDA Parallelism}
For the block diagram with $n$ state variables:
\begin{itemize}[nosep]
  \item Which computations can be parallelized across states?
  \item Which computations must be sequential?
  \item How would you map the $n$ blocks of $f(x)$ and $n \times m$ blocks of $g(x)$ to GPU threads?
\end{itemize}

\begin{answerbox}\end{answerbox}

\question{CUDA Necessity}
Under what conditions would CUDA actually provide a speedup over CPU computation for this problem? (Hint: think about the size of $n$ and $m$, and the number of simulation steps.)

\begin{answerbox}\end{answerbox}

%──────────────────────────────────────────────────
\subsection{Implementation Comparison Framework}
%──────────────────────────────────────────────────

\subsubsection*{Synthesis Question}

\question{Comparison Criteria}
The project requires comparing implementations across languages. Fill in the table:

\medskip
\begin{tabularx}{\textwidth}{l X X}
\toprule
\textbf{Criterion} & \textbf{How to Measure} & \textbf{Expected Winner} \\
\midrule
Execution speed        & & \\[0.8em]
Code complexity        & & \\[0.8em]
Numerical accuracy     & & \\[0.8em]
Ease of visualization  & & \\[0.8em]
Real-time capability   & & \\
\bottomrule
\end{tabularx}

\begin{answerbox}\end{answerbox}

%──────────────────────────────────────────────────
\subsection*{Self-Assessment Checklist}
%──────────────────────────────────────────────────

\begin{itemize}[leftmargin=2em, label=$\square$]
  \item I can identify the MATLAB functions needed for control system simulation
  \item I can write a basic simulation loop in MATLAB pseudocode
  \item I can map MATLAB tools to Python equivalents (NumPy, SciPy, \texttt{python-control})
  \item I can implement one time step of the system model in Python/NumPy
  \item I can explain why C++ and CUDA matter for real-time control
  \item I can identify which computations parallelize on GPU
  \item I can define comparison criteria for cross-language implementation evaluation
\end{itemize}

\begin{notesbox}\end{notesbox}
